% !TeX root = ../papers/nc-p3.tex
\section{三中心商の乗法鉛筆と例外超平面}
\label{sec:three-center-pencil}

\begin{lemma}\label{lem:three-center-pencil}
四周期経路代数で三つの中心四次式を割った商を $C$ とする。
始点 $p$ の長さ $n$ の成分を $C_{p,n}$ と書く。次を仮定する。
\begin{enumerate}
\item $0\le n\le13$ で $\dim C_{p,n}$ は $(1-t^4)^3/(1-t)^4$ の次数 $n$ の係数である。
\item $0\le n<9$ で、$C_{p,n}$ の元のうち右から全ての矢を掛けて零になるものは零だけである。
\end{enumerate}
残る中心四次式の線型結合 $z$ について、$(C/zC)_{p,n}$ の次元が
次数 $0$ から $13$ で $(1+t+t^2+t^3)^4$ の係数に一致することと、
$z_p:C_{p,9}\longrightarrow C_{p,13}$ が同型であることは同値である。
\end{lemma}
\begin{proof}
両空間の次元は64であり、余核が十三次の商である。
核に入る長さ $m<9$ の非零元があれば、第二仮定を反復して右に矢を掛け、
長さ9の非零元を得る。この元も左乗法 $z_p$ の核に入り矛盾する。
従って各 $m\le9$ で単射であり、余核の次元は四つ前の係数との差になる。
逆は十三次の余核の消滅による。
\end{proof}

\begin{proposition}\label{prop:f81-center-pencil}
第\ref{sec:rank-fifteen-fixed-center-obstruction}節の $\mathbb F_{81}$ の表示
$(c,v)=(a^2+2,a^3+a^2+a)$ で、中心 $z_1,z_2,z_5$ を固定する。
第四式を $z_4+s z_3$ とする射影直線全体で、四始点の商が十三次まで所定の次元を持つための
必要十分条件は
\[
       s\notin\{0,-1,-(2a^3+a+2)\}.
\]
これは幾何点について成立し、$s=\infty$ は第四式 $z_3$ を表す。
\end{proposition}
\begin{proof}
\path{computations/verify_rank_fifteen_center_pencil.py} は三式だけで語を簡約する。
四始点のそれぞれで長さ0からの次元は
$1,4,10,20,32,44,54,60,63,64,64,64,64,64$ となる。
各文字次数で四つの右乗法の行列を横に連結し、長さ0から8で行が独立であることを確かめる。
長さ9と13の各文字空間は16次元である。三中心商の同一の語基底において
$z_4,z_3$ の左乗法行列を $M,N$ とすると、全十六ブロックで
\[
 \frac{\det(M+sN)}{\det N}=s^8(s+1)^4(s+2a^3+a+2)^4 .
\]
実際、行列式は高々十六次なので相異なる十七点での有限体消去と補間で全係数が決まり、
最高次係数も非零である。補題\ref{lem:three-center-pencil}を適用すればよい。
三例外で四式を直接課しても、長さ9から13で $s=0$ は $(23,22,22,22,22)$、
他の二値は $(20,16,16,16,16)$ となる。一方、五中心を全て課した商の長さ0から8の
次元は四始点とも $(1,4,10,20,30,36,34,20,0)$ である。従ってこの表示では
元の代数の各始点は五中心上で長さ7以下の155本の経路で生成される。実際、
商基底の持上げを引けば正次数の中心倍が残り、長さに関する帰納法が使える。
\end{proof}

\begin{proposition}\label{prop:five-center-bad-hypersurface}
標数零で五次元の中心四次式空間 $V$ を指定し、折り畳み環 $A$ が $k[V]$ 上有限で
標準Koszul左右AS正則であるとする。四平面 $U\subset V$ を双対射影空間
$\mathbb P(V^*)$ の点で表す。$A/(U)$ が無限次元となる点の集合は、
一つの頂点の隅から得られる次数64の斉次式の零点集合である。
特に有限性を実現する超平面を探す際、任意に固定したアフィン座標で
\[
 (i_1,i_2,i_3,i_4),\qquad i_j\in\mathbb Z_{\ge0},\quad \sum i_j\le64
\]
の有限リストが全てのこのような窓に共通して使える。
\end{proposition}
\begin{proof}
$k[V]$ は換算次数1で生成され、有限性と標準Hilbert関数から次元4である。
一般の四元 $x_1,\ldots,x_4$ を斉次正規化に取り、$z$ で $V$ の基底に延ばす。
$P=k[x_1,\ldots,x_4]$ とする。定理\ref{thm:central-window-cutoff}の必要性の証明により
$A$ は $P$ 自由であり、隅 $D=e_0 A e_0$ の商基底も斉次自由基底に持ち上がる。
経路長 $0,4,8,12$ の基底数は $1,31,31,1$ である。$D$ 上の $z$ の乗法を
この64基底で $M(x)$ と書くと $F(x,u)=\det(uI-M(x))$ は $u$ についてモニックである。
入力と出力の基底次数の差が置換積で相殺するため、換算次数で全次数64の斉次式となる。

$k[X_1,\ldots,X_4,u]$ から $A$ への中心作用を考える。閉点 $(b,d)$ の上での
$D$ のファイバーは $dI-M(b)$ の余核であり、その台は $V(F)$ である。
原点外ではこの台は $A$ の台とも一致する。実際、非零ファイバーで
ある中心四次式が非零スカラーとなる。各長さ4の閉路は頂点0を通るので、
このファイバーでは $1$ が $e_0$ の生成する両側イデアルに入る。
従ってファイバーが非零ならその第0隅も非零である。
超平面 $U$ を割ったときの台はこの台と $U$ の切り出す直線との交わりである
（有限加群に対する中山の補題）。斉次性により、その商が有限次元であることと
直線の射影点で $F$ が消えないことが同値となる。
この記述は線型座標変換で保たれる。

どのアフィンチャートでも $F$ の非斉次化は非零で次数は高々64である。
次数64以下の多項式に対する基底
$\prod_{j=1}^4\binom{t_j}{m_j}$（$\sum m_j\le64$）を取る。
指定した格子での評価行列は全次数順で対角成分1の三角行列である。
よって少なくとも一つの値で $F\ne0$ となる。
\end{proof}

この有限リストの各四平面について、十三次の消滅をまず一始点だけで検査できる。
そこを通過するまで高々三文字なので、その始点で消えれば全始点の長さ16以上が消える。
このとき中心商は有限であり、定理\ref{thm:central-window-cutoff}の二条件でAS性が決まる。
指定した五中心上の有限性が既知なら、全リストで一始点の十三次が非零のときも
命題の対偶と同定理の必要条件により非ASである。五中心上の有限性そのものは、
全五式で割った商の一始点で一つの次数が零になれば従う。

\begin{proposition}\label{prop:five-center-generic-eight}
標数3の二変数単元チャートの始点0で、規格化された固定四列 $(1,2,4,5)$ の
九次商は、全ての幾何点で少なくとも23次元である。
これに対して五列を全て課すと、関数体上の始点0の商次元は次数0から8で
\[
              1,4,10,20,30,36,34,20,0
\]
となる。
\end{proposition}
\begin{proof}
\path{computations/verify_rank_fifteen_full_center_cut.py} は
第\ref{sec:rank-fifteen-bivariate-sixth-cut}節の語を末尾へ延ばす再帰を
$\mathbb F_3(K,H)$ 上で実行する。四列の九次の文字別次元は $(6,6,6,5)$ である。
代入可能性には、最後まで同じ消去主元を取れることを仮定しない。
無切断の五次までの語基底と四次中心の規格化にはもとのチャートの単元だけが必要である。
長さ9の全語に二次関係と規格化四列を挿入した表示に戻ると、
余核の次元はどの特殊化でも一般点以上である。平方根による語の基底変換は
単元の体拡大上で実行でき、次元を変えない。この議論で得る下界は標数3に関するものである。
五列の場合、六次から八次までの四文字の次元は
$(10,10,10,4),(5,5,5,5),(0,0,0,0)$ である。
\end{proof}

\begin{theorem}\label{thm:ordinary-fifteen-finite-over-quartics}
正の固定曲線の標数零の全通常階数15窓で、折り畳み環 $A$ は四次中心の生成する部分環上で有限である。
四次中心全体で割った商は各始点で長さ13において零となる。
この族では次の手順が、非正則の場合も含め、各入力の標準Koszul左右AS正則性を決定する。
四次中心の基底を長さ5の交換子で求め、その四本の一般線型結合で割る。
係数を独立な不定元として、四始点・長さ0から13で商の次元が
$(1+t+t^2+t^3)^4$ の係数以下かを検査する。この条件は必要十分である。
\end{theorem}

\begin{proof}
四つの中心四次式上ですでに自由性が証明された窓には有限性がある。
第\ref{sec:rank-fifteen-quadratic-parabolic-pencils}節の帰着と
定理\ref{thm:rank-fifteen-all-quadratic-verticals},
\ref{thm:rank-fifteen-prime-parabolic-pencils}により、残るのは3上の付値で
$X=\bar x\notin\mathbb F_9$, $Y=\bar y\ne0,(1-X)^2$ の単元対である。
$X=c^2$, $B=Y/X^3=(v^3-v)^2$ とする。
前命題の五中心の八次までの主元を集め、開集合の単元を除いた
$G(v,c)$ の二次数は $(177,140)$ である。
全十二表示で $G=0$ となる必要十分条件は、自由基底
$c^i v^j$ ($0\le i<2$, $0\le j<6$) における十二係数が全て零となることである。
ここでは $XB\ne0$ により各ファイバーが被約である。

検証コードは十二係数の最大公約元が1であることを確認する。
第\ref{sec:rank-fifteen-bivariate-seventh-cut}節の手順で各係数から
$X, X-1,B,X^3B-(1-X)^2$ の因子を除き、各 $i$ で $(j,j+1)$
($j=0,2,4$) の組から $B$ を消去する。$X, X-1$ の因子を除いた
六つの終結式のモニック最大公約元は
\[
 (X+1)^{10}(X^2+1)^2 (X^3+2X+2)(X^5+2X^4+2X+2)
\]
である。後の二既約因子の剰余体で、除去前の十二係数を $B$ の多項式として
最大公約元を取ると、その全幾何根はそれぞれ
$B=X^2+2X+1$ と $B=2X^3+X^2$ の一点ずつである。
これらの二軌道では元の四窓を組み立て、四中心商の標準次元を
四始点・長さ13まで直接照合する。使用した体の絶対拡大次数は順に6と5である。
従ってここでは混標数への既出の四切断の持上げを使える。

その他の点では少なくとも一表示で $G\ne0$ なので、五中心商は始点0の長さ8で零である。
第\ref{sec:rank-fifteen-truncated-central-freeness}節と同じくモニックな座標方程式で
この表示を持ち上げる。長さ5までの経路基底が持ち上がり、中心方程式は剰余で核次元5、
標数零でも一般窓からの下界5を持つため、五つの自由列を全て持ち上げられる。
語の有限表示の余核次元は分数体で増えない。他の始点でも長さ11の経路は
始点0からの長さ8の区間を含むので、五中心商は有限次元である。
中心倍を除いて長さに関して帰納すれば、元の経路代数が五中心上で有限となる。
ここで目標と同じ $(x,y)$ を持つ表示の選択は十分である。
定理\ref{thm:rank-fifteen-rational-quotient}の証明の巡回移動、語反転、基底変更は
四次中心上での有限性を保つ。この理由で、固定した始点の商次元をそのまま運ぶときに
必要だった符号の制約を課さずに十二表示を使える。
四中心上の自由性を使った場合には四中心商の長さ13が零であり、
残りの場合は上の五中心商の長さ11が零なので、全四次中心で割った商はどちらも長さ13で零になる。
体拡大後の有限性も降下する。実際、四次中心全体は有限行列の核で体拡大と可換であり、
その商の高次消滅は忠実平坦な拡大から降りる。

正則な入力では有限性と斉次Noether正規化により、四つの一般四次中心が
有限次元商を与え、定理\ref{thm:central-window-cutoff}の必要条件から所定の係数に一致する。
逆に、正の固定曲線のこの族は第二補完の閉包の通常窓であり、当該上界が成立すれば
命題\ref{prop:klein-central-strata}の自由性・正則性の十分条件が適用できる。
これで両方向が従う。添字保持同型は定理\ref{thm:rank-fifteen-global-invariants}の
$(x,p,q)$ で判定できる。この有限手順から、全通常点で正則性の答が一律に肯定となることまでは結論しない。
\end{proof}
